% GENERATED FILE. Edit canonical lessons, then run npm run generate:books.
% canonical-source-hash: fnv1a64:192795b3db00fea7
% canonical-lessons: FR-C54-hopital, FR-C54-maladie, FR-C54-soif, FR-C54-fatigue, FR-C54-malade

\chapter{Hospital, Illness, Thirst, Tiredness, Ill}
\label{ch:fr-hospital-illness-thirst-tiredness-ill}
\hlchaptermodality{54}

\begin{chapteropening}
\textbf{By the end of this chapter:} \emph{I can name a hospital, an illness, thirst, tiredness, ill.}

Complete the last lesson of chapter 54: I can name a hospital, an illness, thirst, tiredness, ill.
\end{chapteropening}

\section[l'hôpital]{l'hôpital — a hospital}

\label{lesson:FR-C54-hopital}



\begin{quote}
Before the new one: say the French for a medicine, then the French for a doctor.
\end{quote}

\subsection*{You'll want to know: l'hôpital}
\textbf{l'hôpital} — \textquotedblleft{}a hospital\textquotedblright{}.

\begin{cousinweb}
\textbf{hôpital} $\leftarrow$ Old French \emph{hospital} $\leftarrow$ Medieval Latin \emph{hospitale}, \textquotedblleft{}a guest house\textquotedblright{}, from Latin \emph{hospes}, \textquotedblleft{}guest\textquotedblright{} — a learned borrowing. (EtymDB 2.1) The circumflex marks the \emph{s} that French later stopped writing. The same Latin word, worn down in everyday speech, also gave French its word for \textquotedblleft{}hotel\textquotedblright{}.
\end{cousinweb}

\begin{grammarlens}[title={What you've built}]
One more word for this chapter.
\end{grammarlens}

\subsection*{Your turn}
\begin{itemize}
\raggedright
  \item \emph{Say it:} \emph{l'hôpital}
  \item \emph{Say it:} \emph{l'hôpital}, once more
  \item \emph{Say it:} say \emph{l'hôpital}
  \item \emph{Recall:} say the French for a medicine, then the French for a doctor, then say \emph{l'hôpital} again
\end{itemize}

\begin{tcolorbox}[breakable,colback=teal!4,colframe=teal!35!black,title={Before you move on}]
What does l'hôpital mean? (\textquotedblleft{}A hospital\textquotedblright{}.) Say it once more.
\end{tcolorbox}
\section[la maladie]{la maladie — an illness}

\label{lesson:FR-C54-maladie}



\begin{quote}
Before the new one: say the French for a doctor, then the French for a hospital.
\end{quote}

\subsection*{You'll want to know: la maladie}
\textbf{la maladie} — \textquotedblleft{}an illness\textquotedblright{}.

\begin{cousinweb}
\textbf{maladie} — built in Old French from the adjective for \textquotedblleft{}sick\textquotedblright{}, with the ending \emph{-ie}. That adjective goes back to Latin \emph{male habitus}, \textquotedblleft{}in poor condition\textquotedblright{}. (EtymDB 2.1)
\end{cousinweb}

\begin{grammarlens}[title={What you've built}]
One more word for this chapter.
\end{grammarlens}

\subsection*{Your turn}
\begin{itemize}
\raggedright
  \item \emph{Say it:} \emph{la maladie}
  \item \emph{Say it:} \emph{la maladie}, once more
  \item \emph{Say it:} say \emph{la maladie}
  \item \emph{Recall:} say the French for a doctor, then the French for a hospital, then say \emph{la maladie} again
\end{itemize}

\begin{tcolorbox}[breakable,colback=teal!4,colframe=teal!35!black,title={Before you move on}]
What does la maladie mean? (\textquotedblleft{}An illness\textquotedblright{}.) Say it once more.
\end{tcolorbox}
\section[la soif]{la soif — thirst}

\label{lesson:FR-C54-soif}



\begin{quote}
Before the new one: say the French for a hospital, then the French for an illness.
\end{quote}

\subsection*{You'll want to know: la soif}
\textbf{la soif} — \textquotedblleft{}thirst\textquotedblright{}. \textbf{j'ai soif}: I am thirsty, with \textbf{avoir}.

\begin{cousinweb}
\textbf{soif} $\leftarrow$ Old French \emph{soif} $\leftarrow$ Latin \emph{sitis}, \textquotedblleft{}thirst\textquotedblright{} — inherited. (EtymDB 2.1)
\end{cousinweb}

\begin{grammarlens}[title={What you've built}]
One more word for this chapter.
\end{grammarlens}

\subsection*{Your turn}
\begin{itemize}
\raggedright
  \item \emph{Say it:} \emph{la soif}
  \item \emph{Say it:} \emph{la soif}, once more
  \item \emph{Say it:} say \emph{la soif}
  \item \emph{Recall:} say the French for a hospital, then the French for an illness, then say \emph{la soif} again
\end{itemize}

\begin{tcolorbox}[breakable,colback=teal!4,colframe=teal!35!black,title={Before you move on}]
What does la soif mean? (\textquotedblleft{}Thirst\textquotedblright{}.) Say it once more.
\end{tcolorbox}
\section[la fatigue]{la fatigue — tiredness}

\label{lesson:FR-C54-fatigue}



\begin{quote}
Before the new one: say the French for an illness, then the French for thirst.
\end{quote}

\subsection*{You'll want to know: la fatigue}
\textbf{la fatigue} — \textquotedblleft{}tiredness\textquotedblright{}.

\begin{cousinweb}
\textbf{fatigue} — formed from the verb for \textquotedblleft{}to tire\textquotedblright{}, a learned borrowing of Latin \emph{fatigare}, \textquotedblleft{}to weary\textquotedblright{}. (EtymDB 2.1) The kept \emph{g} shows the book word: in everyday speech it would have dropped out.
\end{cousinweb}

\begin{grammarlens}[title={What you've built}]
One more word for this chapter.
\end{grammarlens}

\subsection*{Your turn}
\begin{itemize}
\raggedright
  \item \emph{Say it:} \emph{la fatigue}
  \item \emph{Say it:} \emph{la fatigue}, once more
  \item \emph{Say it:} say \emph{la fatigue}
  \item \emph{Recall:} say the French for an illness, then the French for thirst, then say \emph{la fatigue} again
\end{itemize}

\begin{tcolorbox}[breakable,colback=teal!4,colframe=teal!35!black,title={Before you move on}]
What does la fatigue mean? (\textquotedblleft{}Tiredness\textquotedblright{}.) Say it once more.
\end{tcolorbox}
\section[malade]{malade — ill, sick}

\label{lesson:FR-C54-malade}



\begin{quote}
Before the new one: say the French for thirst, then the French for tiredness.
\end{quote}

\subsection*{You'll want to know: malade}
\textbf{malade} — \textquotedblleft{}ill, sick\textquotedblright{}.

\begin{cousinweb}
\textbf{malade} $\leftarrow$ Old French \emph{malade} $\leftarrow$ Latin \emph{male habitus}, \textquotedblleft{}in poor condition\textquotedblright{} — inherited. (EtymDB 2.1)
\end{cousinweb}

\begin{grammarlens}[title={What you've built}]
One more word for this chapter.
\end{grammarlens}

\subsection*{Your turn}
\begin{itemize}
\raggedright
  \item \emph{Say it:} \emph{malade}
  \item \emph{Say it:} \emph{malade}, once more
  \item \emph{Say it:} say \emph{malade}
  \item \emph{Recall:} say the French for thirst, then the French for tiredness, then say \emph{malade} again
\end{itemize}

\begin{tcolorbox}[breakable,colback=teal!4,colframe=teal!35!black,title={Before you move on}]
What does malade mean? (\textquotedblleft{}Ill, sick\textquotedblright{}.) Say it once more.
\end{tcolorbox}
